\subsection{Integral induction in wild ramification}\label{sec:second-two-wild}
The object of this subsection is the following theorem, which completes the computation of $K_2$ of division algebras over two-dimensional local fields.

\begin{theorem}\label{thm:second-two-vanishing}
Let $F$ be a two-dimensional local field.
For every central division algebra $D$ over $F$, one has, integrally,
\[
 \SK_2(D)=0,\qquad \Nrd_2\colon K_2(D)\xrightarrow{\ \sim\ }K_2(F).
\]
\end{theorem}

The proof of Theorem~\ref{thm:second-two-vanishing} is given at the end of this subsection, after that of Theorem~\ref{thm:second-two-wild}.

\begin{theorem}\label{thm:second-two-wild}
Let $k$ be a nonarchimedean local field of characteristic $p>0$, and let $F$ be a complete discretely valued field with first residue field $k$.
The field $F$ may have characteristic $p$ or zero.
For every $p$-primary central division algebra $D/F$, one has $\SK_2(D)=0$.
\end{theorem}

We outline the proof of Theorem~\ref{thm:second-two-wild}.
It is an induction on the minimal unramified index.
For a finite Galois unramified extension $F'/F$ realizing the minimal index, the kernel of restriction on $\SK_2$ is contained in the corestriction of $\SK_2(D')$, where $D'$ is the division algebra in the Brauer class of $D_{F'}$; hence it suffices to prove $\SK_2(D')=0$, and we may assume that $\deg D=\ind(D_{F^{\mathrm{nr}}})=p^q$.
If $q=0$, then $D$ is split.

If $q>0$, after adjoining $\mu_p$ if necessary and passing to a further finite Galois unramified extension, the cyclic cover lemma gives a cyclic extension $T/F$ of degree $p$ inside $D$ whose centralizer $B$ has minimal unramified index $p^{q-1}$.
This is where the induction hypothesis enters: applied to $B$, it gives $\SK_2(B)=0$, and together with the isomorphism of the reduced norm on splitting transfers it gives $K_2(B)=I_T(B)$.
This handles the generic images of the corner components of the orbits of the relative cover, and the lemma on orbit images handles the other components, so the generic image of the whole $\pi_2$ of the orbits lies in $I_F(D)$.

Finally, for $\xi\in\SK_2(D)$, we lift $\xi$ to an integral class with zero residue, lift its trace through the surjection on $\pi_2$ of degree-zero $\mathrm{TC}$, and remove the residue component.
At a fixed finite truncation we represent the remaining relative orbit class by an integral class and, in the mixed-characteristic homogeneous-unit type, add back an integral representative of the class subtracted with the residue.
The generic image of the resulting class lies in $I_F(D)$; its difference from the original lift vanishes at that truncation, so the generic image of the difference lies in $I_F(D)$ by congruence separation (Lemma~\ref{lem:app-congruence-separation}).
Hence $\xi\in I_F(D)$, and since $\Nrd_2$ is injective on $I_F(D)$, we conclude that $\xi=0$.

\subsubsection{Unramified descent and minimal index}
We first reduce Theorem~\ref{thm:second-two-wild} to division algebras whose index does not drop over the maximal unramified extension.
Recall from Definition~\ref{def:paper-transfer} the subgroup $I_F(D)=\sum_E\theta_E K_2(E)$, with $E/F$ ranging over finite splitting fields.

\begin{proposition}
\label{prop:review-splitting-transfer-norm-iso}
For a two-dimensional local field $F$ and any central simple algebra $D/F$, the map $\Nrd_2\colon I_F(D)\xrightarrow{\sim}K_2(F)$ is an isomorphism.
\end{proposition}
\begin{proof}
If $\ell\ne\operatorname{char}F$, then $\cd_\ell F\le3$, so Corollary~\ref{cor:review-small-cd-splitting-transfer} applies.
If $\operatorname{char}F=p$, then $[F:F^p]=p^2$, and Lemma~\ref{lem:review-p-rank-two-frobenius-norm} handles the remaining $p$-component.
Since $\Nrd_2\colon I_F(D)\to K_2(F)$ is an isomorphism after localization at every prime, the integral map is also an isomorphism.
\end{proof}

\begin{lemma}
\label{lem:review-minimal-order-k1-support}
Let $F$ be a two-dimensional local field whose first residue field $k$ has characteristic $p$, and let $D/F$ be a central division algebra of degree $n=p^q$.
Assume $\ind D_{F^{\mathrm{nr}}}=n$, and set $P=\mathcal O_D$ and $\Delta=P/J(P)$.
Then the ramification index is $n$, the residue field is $\Delta=k^{1/n}$, and $\per D_{F^{\mathrm{nr}}}=n$.
Moreover, the sequence
\[
 0\longrightarrow\Delta^\times/k^\times
 \longrightarrow K_1(P)\longrightarrow
 \Nrd_1(P^\times)\longrightarrow0
\]
is exact.
For a finite unramified extension $F'/F$, the map $K_1(P)\to K_1(\mathcal O_{D_{F'}})$ is injective, as is the map to the direct limit over all finite unramified extensions.
\end{lemma}
\begin{proof}
At every finite unramified stage, which is again a two-dimensional local field, the index stays $n$, and so does the period by the equality of period and index over every finite extension \cite[Corollary 5.1, Proposition 5.2]{Chipchakov2022}.
A smaller positive multiple of the Brauer class vanishing over $F^{\mathrm{nr}}$ would vanish at a finite stage by descent of its splitting data, so $\per D_{F^{\mathrm{nr}}}=n$.
If the residue center had a nontrivial separable part, or the residue algebra had a noncommutative division-algebra part, a finite separable residue extension splitting that part would produce a nontrivial idempotent.
Henselianity lifts that idempotent to the order, contradicting the fact that the generic fiber is a division algebra.
Thus $\Delta/k$ is purely inseparable.
Since $[k:k^p]=p$, the extension is simple; lifting a primitive element gives $f=[\Delta:k]\le n$.
Lifting a generator of the value group similarly gives $e\le n$ for the ramification index.
The defectless dimension formula $n^2=ef$ now yields $e=f=n$, and hence $\Delta=k^{1/n}$ \cite{WadsworthValuation}.

Let $i_*:K_1(\Delta)\to K_1(P)$ be the support map.
For $x\in K_2(D)$, put $c=\partial_Dx$.
For a separable maximal subfield $E\subset D$ and the split Morita equivalence $\mu_E$, put $\rho_E=\mu_E^{-1}\res_{E/F}\colon K_2(D)\to K_2(E)$.
Then $\theta_E\circ\rho_E=n$, $N_{E/F}\circ\rho_E=N_{D/F}$, and $\Nrd_2\theta_E=N_{E/F}$ \cite[Proposition 6.6.1(1)]{KahnLevine2010}, so $N_{D/F}=\Nrd_2\theta_E\rho_E=n\Nrd_2$.
The norm $N_{D/F}=n\Nrd_2$ satisfies $\partial_FN_{D/F}=N_{\Delta/k}\partial_D$, and $N_{\Delta/k}(c)=c^n$ because $\Delta/k$ is purely inseparable of degree $n$, so
\[
 c^n=(\partial_F\Nrd_2x)^n.
\]
Since $\Delta^\times$ has no $p$-torsion, it follows that $c=\partial_F\Nrd_2x\in k^\times$.
The inclusion $k^\times\subseteq\partial_DK_2(D)$ follows from Theorem~\ref{thm:second-two-surj} and surjectivity of the boundary for the commutative field, so $\ker i_*=k^\times$.
Since $\SK_1(D)=0$ by Corollary~\ref{cor:review-two-local-sk1-zero}, we have $\ker[K_1(P)\to\Nrd_1(P^\times)]=\im i_*$, which proves the displayed exact sequence.

For a finite unramified extension with residue extension $k'/k$, the new residue field is $\Delta'=\Delta k'$.
The boundary comparison alone gives $\ker i_*'\subseteq(k')^\times$.
If $z\in K_1(P)$ becomes zero after base change, then $\Nrd_1(z)=1$, so $z=i_*c$ for some $c\in\Delta^\times/k^\times$.
Vanishing after base change forces $c\in(k')^\times$.
The linear disjointness of purely inseparable and separable extensions gives $\Delta\cap k'=k$, hence $z=0$.
Injectivity of the map to the direct limit follows from continuity of $K$-theory.
\end{proof}

\begin{lemma}
\label{lem:review-geometric-bockstein-detection}
In the situation of the preceding lemma, let $P^{\mathrm{nr}}$ be the union of all finite unramified base changes of $P$, and let $\pi$ be a uniformizer of $\mathcal O_F$.
There is $j_0$ such that, for all $a\ge1$ and $j\ge j_0$,
\[
 K_1(P)[p^a]\longrightarrow K_1(P^{\mathrm{nr}}/\pi^jP^{\mathrm{nr}})
\]
is injective.
Consequently, if a class in $K_2(P;\Z/p^a)$ vanishes on this geometric truncation, its Bockstein is zero.
\end{lemma}
\begin{proof}
We use the injectivity in Lemma~\ref{lem:app-deep-unit-norm}, uniform over unramified extensions.
The hypothesis $\ind D_{F^{\mathrm{nr}}}=n$ says that every $D_L$ is a division algebra of the same degree and that $P_L$ is obtained by base change from the same integral lattice.
Thus there is $j_1$ such that, for every finite unramified $L/F$ and $j\ge j_1$, the map $\Nrd_1$ is injective on the image of $1+\pi^jP_L$ in $K_1(P_L)$.

In mixed characteristic put $e=v_F(p)$.
For $j>e/(p-1)$, the group $1+\pi^j\mathcal O_L$ contains no nontrivial root of unity of $p$-power order.
Indeed, if $\zeta\ne1$ has order $p^r$, then $v_L(\zeta-1)=e/(p^{r-1}(p-1))$, and $e$ is constant under unramified extension.
In equal characteristic the multiplicative group has no $p$-torsion.
Choosing $j_0$ sufficiently large therefore makes the $\pi^j$-congruence kernel in every $K_1(P_L)$ free of $p$-primary torsion.
Here lifting general linear groups shows that this congruence kernel equals $\im(1+\pi^jP_L\to K_1(P_L))$.

If $z\in K_1(P)[p^a]$ vanishes on the geometric truncation, it already vanishes at a finite stage $L/F$.
Then $z_L$ is a $p$-primary torsion element of the preceding congruence kernel, hence is zero.
The unramified injectivity in Lemma~\ref{lem:review-minimal-order-k1-support} gives $z=0$.
The last assertion follows by naturality of the Bockstein in the Moore coefficient sequence.
\end{proof}

\begin{lemma}
\label{lem:review-geometric-moore-integral-lift}
Let $F$ be a two-dimensional local field whose first residue field $k$ has characteristic $p$, and let $B/F$ be a central division algebra with $n=p^q=\deg B=\ind B_{F^{\mathrm{nr}}}$.
Let $R=\mathcal O_F$, $C=\mathcal O_B$, $\Delta=C/J(C)$, and let $\pi$ be a uniformizer of $R$.
Let $Y_C=K(C,\Delta)$ and $Y_s=K(C/\pi^sC,\Delta)$.
For every $a,j\ge1$, there is $b\ge j$ such that, for every geometrically zero class in $\pi_2(Y_b/p^a)$, the image of that class in $Y_j/p^a$ is represented by an integral element of $\pi_2Y_C$.
\end{lemma}
\begin{proof}
We lift a generator of $\Delta/k$ to $w\in C^\times$.
Since $[\Delta:k]=\deg B=n$ by Lemma~\ref{lem:review-minimal-order-k1-support}, the residue degree of $E=F(w)$ is at least $n$.
The bound on the degree of a maximal subfield gives $[E:F]=n$ and ramification index one.
The finite free $R$-algebra $R[w]$ has residue field $\Delta$ and maximal ideal $\pi R[w]$, so it is a discrete valuation ring, and hence $R[w]=\mathcal O_E\subset C$.
Since both entries of a Matsumoto symbol in $K_2(\Delta)$ lift to units of this commutative order, the map $K_2(C)\to K_2(\Delta)$ is surjective, and the same lifts work in every truncation.
Thus
\[
\begin{gathered}
 \pi_1Y_C=M,\qquad M=\ker[K_1(C)\to\Delta^\times],\\
 \pi_1Y_s=M/\mathcal U_s(C),\qquad \mathcal U_s(C)=\im(1+\pi^sC\to K_1(C)).
\end{gathered}
\]
By Corollary~\ref{cor:app-integral-cut-stable}, for sufficiently deep $r\ge j$ the image of $\pi_2Y_r$ in $\pi_2Y_j$ equals the image of $\pi_2Y_C$.
Increase $r$ so that $r\ge j_0$ for the integer $j_0$ of Lemma~\ref{lem:review-geometric-bockstein-detection} and so that, by Lemma~\ref{lem:app-deep-unit-norm}, the reduced norm is injective on $\mathcal U_r(C)$.

For a sufficiently deep $b$, write the Bockstein of the given class as $\beta_b\in(M/\mathcal U_b(C))[p^a]$ and lift it to $\beta\in M$.
Then $p^a\beta\in \mathcal U_b(C)$.
Put $t=\Nrd_1(\beta)$, so $t^{p^a}\in1+\pi^bR$.
In mixed characteristic, if $b$ is sufficiently large, Hensel's lemma provides a $p^a$-th root $\eta$ of $t^{p^a}$ in $1+\pi^{r+c}R$, where $c$ satisfies $1+\pi^{r+c}R\subset\Nrd_1(\mathcal U_r(C))$.
In equal characteristic, $v_F(t^{p^a}-1)=p^a v_F(t-1)$, so $b\ge p^a(r+c)$ allows us to take $\eta=t\in1+\pi^{r+c}R$.
We choose $\gamma\in \mathcal U_r(C)$ with $\Nrd_1(\gamma)=\eta$.
The element $p^a(\beta-\gamma)$ belongs to $\mathcal U_r(C)$ and has reduced norm one, so it is zero.
Thus $\beta_0=\beta-\gamma\in M[p^a]$ represents the image of $\beta_b$ at depth $r$.

Since the given Moore class is geometrically zero, $\beta_0$ vanishes on the geometric truncation at depth $r$.
Lemma~\ref{lem:review-geometric-bockstein-detection} gives $\beta_0=0$.
The image of the given class in $Y_r/p^a$ therefore has zero Bockstein and comes from some $u_r\in\pi_2Y_r$.
By the choice of $r$ via Corollary~\ref{cor:app-integral-cut-stable}, the image of $u_r$ in $Y_j$ is represented by an element of $\pi_2Y_C$.
\end{proof}

Let $F'/F$ be a finite Galois unramified extension of a two-dimensional local field $F$, let $D/F$ be a central division algebra, and let $D'$ be the division algebra in the Brauer class of $D_{F'}$.
Put $R=\mathcal O_F$, $R'=\mathcal O_{F'}$, $P=\mathcal O_D$ and $C=\mathcal O_{D'}$.
The order $P\otimes_RR'$ is hereditary.
Indeed, finite \'etale base change is separable, so every base-changed module is a direct summand of the module obtained by restricting and inducing it again; the bound of one on projective dimension is preserved.
After rectangular diagonal blocks are removed by Morita equivalence, the lattice-chain classification of hereditary orders therefore identifies $P\otimes_RR'$ with an Iwahori order with corner $C$ \cite[Chapter 9]{ReinerMaximalOrders}.

\begin{lemma}
\label{lem:review-corner-h-one-transfer}
Let $F$ be a two-dimensional local field whose first residue field $k$ has characteristic $p$, and let $D/F$ be a central division algebra.
Let $R=\mathcal O_F$, $P=\mathcal O_D$, and let $\pi$ be a uniformizer of $R$.
Let $a\ge1$.
Let $F'/F$ be a finite Galois unramified extension, put $R'=\mathcal O_{F'}$, and let $k'/k$ be the residue extension.
Let $D'$ be the division algebra in the Brauer class of $D_{F'}$, and put $C=\mathcal O_{D'}$.
After removing rectangular diagonal blocks by Morita equivalence, we write $P\otimes_RR'$ as an Iwahori order $H$ with corner $C$.
Let
\[
 \mathcal X_s=K(P/\pi^sP,P/J(P))/p^a,\qquad
 \mathcal Y_s=K(C/\pi^sC,C/J(C))/p^a.
\]
These may also be regarded as sheaves on the finite \'etale site, as in Lemma~\ref{lem:app-nilpotent-etale-descent}.
Suppose that classes $y_s\in\pi_2\mathcal X_s(k)$ at this fixed coefficient $p^a$ are compatible for all $s$ and have zero image at the geometric stalk.
For every $j$ and every prescribed $b_0\ge j$, there are $b\ge b_0$ and a geometrically zero $z_b\in\pi_2\mathcal Y_b(k')$ such that $y_j$ is the image, at depth $j$, of the transfer of $z_b$ from the corner.
\end{lemma}
\begin{proof}
Put $V_s=\pi_3(\mathcal X_s)_{\bar k}$ and $W_s=\pi_3(\mathcal Y_s)_{\bar k}$.
They are discrete modules for $G_k$ and $G_{k'}$, respectively.
Corner induction and finite \'etale norms give a coefficient map
\[
 c_{b,j}\colon\operatorname{Ind}_{G_{k'}}^{G_k}W_b\longrightarrow V_j.
\]
At the geometric stalk, it is the sum of the corner maps indexed by the embeddings $k'\hookrightarrow\bar k$.
Denote its image by $L_{b,j}\subset V_j$.

By Lemma~\ref{lem:app-iwahori-derived-cut}, the transition maps on the relative corner cofibers are nullhomotopic through a factorization via the Koszul algebra defined over $k'$, so the nullhomotopies respect $G_{k'}$.
For fixed $b\ge b_0$, we may therefore take $r$ sufficiently deep that
\[
 V_r|_{G_{k'}}\longrightarrow W_b\longrightarrow V_j|_{G_{k'}}
\]
equals $V_r\to V_j$.
Indeed, if the composite from $\mathcal X_r$ to the corner cofiber of $\mathcal X_b$ is nullhomotopic, lift it to $\mathcal Y_b$ and then pass to depth $j$.
Thus $\im(V_r\to V_j)\subseteq L_{b,j}$.

By Lemma~\ref{lem:app-nilpotent-etale-descent}, $y_r$ comes from some $h_r\in H^1(k,V_r)$.
The image of $h_r$ in $H^1(k,V_j)$ comes from $H^1(k,L_{b,j})$.
The map $\operatorname{Ind}W_b\to L_{b,j}$ is surjective by definition; its kernel is a discrete $p$-primary module, so $H^2(k,\ker(\operatorname{Ind}W_b\to L_{b,j}))=0$.
Consequently,
\[
 H^1(k,\operatorname{Ind}_{G_{k'}}^{G_k}W_b)
 \twoheadrightarrow H^1(k,L_{b,j}).
\]
By Shapiro's lemma, we have $H^1(k,\operatorname{Ind}_{G_{k'}}^{G_k}W_b)\simeq H^1(k',W_b)$, which gives a geometrically zero $z_b\in\pi_2\mathcal Y_b(k')$ through the edge map of Lemma~\ref{lem:app-nilpotent-etale-descent}.
At the geometric stalk, transfer is the sum over the preceding embeddings.
Compatibility with the descent edge map therefore identifies the transfer of $z_b$ at depth $j$ with $y_j$.
\end{proof}

\begin{lemma}\label{lem:second-two-inertial-reduction}
Let $F$ and $D$ be as in Lemma~\ref{lem:review-corner-h-one-transfer}.
Let $F'/F$ be a finite Galois unramified extension with $\ind D_{F'}=n$, where $n=\ind D_{F^{\mathrm{nr}}}$, and let $D'$ be the division algebra in the Brauer class of $D_{F'}$.
Then
\[
 \ker[\SK_2(D)\longrightarrow\SK_2(D')]
 \ \subseteq\ \cor_{F'/F}\SK_2(D').
\]
We do not assume that $[F':F]$ is prime to $p$.
\end{lemma}
\begin{proof}
Put $R=\mathcal O_F$, $P=\mathcal O_D$, $C=\mathcal O_{D'}$, and $R'=\mathcal O_{F'}$.
As explained before Lemma~\ref{lem:review-corner-h-one-transfer}, after rectangular diagonal blocks are removed by Morita equivalence, $P\otimes_RR'$ is an Iwahori order with corner $C$.
An element $\xi$ of $\ker[\SK_2(D)\to\SK_2(D')]$ is $p$-primary torsion, since the prime-to-$p$ parts of $\SK_2$ vanish: for each prime $\ell\ne p$, split the other primary components of $[D]$ over a finite separable extension of degree prime to $\ell$, which detects the $\ell$-primary part by restriction and transfer, and apply Proposition~\ref{prop:second-residue-zero-ell}(2) to the division algebra representing the result.
The residue algebra $P/J(P)$ is a division algebra whose center $l$ is a nonarchimedean local field, so Theorem~\ref{thm:paper-local-norm} gives reduced-norm isomorphisms $K_1(P/J(P))\simeq l^\times$ and $K_2(P/J(P))\simeq K_2(l)$.
The localization boundary of $\xi$ vanishes since $l^\times$ has no $p$-torsion, so $\xi$ lifts to $K_2(P)$.
The support image is $p$-divisible, since $K_2(l)$ is uniquely $p$-divisible \cite{GeisserLevine}.
Adjusting the lift by a support class therefore makes a $p$-power multiple of the adjusted lift zero, so the integral lift itself may be chosen $p$-primary torsion.
Since residue $K_2$ has no $p$-torsion, the adjusted lift has zero reduction, so we choose a relative integral lift $y$ of $\xi$.
Let $\pi$ be a uniformizer of $R$, and put $\Delta'=C/J(C)$.
Let $y'$ be the image of $y$ in $K_2(C)$ after base change and the corner equivalence (Lemma~\ref{lem:app-iwahori-derived-cut}).
Since $\xi|_{D'}=0$, the localization sequence gives $y'\in i_*K_2(\Delta')$, and the same holds after every further finite unramified base change.
Every finite separable extension $\Delta''$ of $\Delta'$ is a nonarchimedean local field of characteristic $p$ with $[\Delta'':\Delta''^p]=p$, so Geisser--Levine gives $K_2(\Delta'';\Z/p^a)=K_3(\Delta'';\Z/p^a)=0$.
The first vanishing kills the image of the closed-support classes in $K_2(C;\Z/p^a)$, and the second makes $\pi_2$ of the relative group inject into the absolute one, so the relative Moore class also vanishes.
Hence, for each fixed $a$, the compatible family has zero image at the geometric stalk, as required in Lemma~\ref{lem:review-corner-h-one-transfer}.

Fix a finite separable splitting field $L_0/F$, and choose a depth $j$ at which Lemma~\ref{lem:app-congruence-separation} holds for $L_0$.
We choose $p^a$ annihilating $K_2(P/\pi^jP,P/J(P))$, which exists because this relative group is $p$-primary of bounded exponent (Lemma~\ref{lem:app-nilpotent-trace}).
Let $b_0$ be a depth at which Lemma~\ref{lem:review-geometric-moore-integral-lift} can be applied to the corner $C$, and apply Lemma~\ref{lem:review-corner-h-one-transfer}.
The lemma applies because $\deg D'=n$ is a power of $p$: the residue field of $F^{\mathrm{nr}}$ is separably closed and $\cd_\ell(F^{\mathrm{nr}})\le1$ for $\ell\ne p$, so $\Br(F^{\mathrm{nr}})$ is $p$-primary; moreover, $D'$ has minimal unramified index $n$, since $(F')^{\mathrm{nr}}=F^{\mathrm{nr}}$.
At depth $j$, the image of the geometrically zero class $z_b$ is represented by an integral $z\in K_2(C,C/J(C))$.
Thus $y$ and the transfer of $z$ from the corner agree in $\pi_2\mathcal X_j$, with $\mathcal X_j$ as in Lemma~\ref{lem:review-corner-h-one-transfer}.
Since $p^a$ kills $K_2(P/\pi^jP,P/J(P))$, the Moore sequence injects this integral group into $\pi_2\mathcal X_j$, so the equality holds integrally at depth $j$.
By Lemma~\ref{lem:app-congruence-separation}, we have $\xi-\cor_{F'/F}(z|_{D'})\in I_F(D)$.
By Proposition~\ref{prop:review-splitting-transfer-norm-iso}, $K_2(D')=I_{F'}(D')\oplus\SK_2(D')$.
Write $z|_{D'}=z_I+z_0$ with $z_I\in I_{F'}(D')$ and $z_0\in\SK_2(D')$.
By transitivity of transfer, $\cor_{F'/F}I_{F'}(D')\subseteq I_F(D)$, so $\xi-\cor_{F'/F}(z_0)\in I_F(D)$.
Since $\Nrd_2\circ\cor_{F'/F}=N_{F'/F}\circ\Nrd_2$, the same element lies in $\SK_2(D)$.
As $I_F(D)\cap\SK_2(D)=0$ by the injectivity of $\Nrd_2$ on $I_F(D)$ (Proposition~\ref{prop:review-splitting-transfer-norm-iso}), we obtain $\xi=\cor_{F'/F}(z_0)$, and hence
\[
 \xi\in\cor_{F'/F}\SK_2(D').
\]
\end{proof}

By Lemma~\ref{lem:second-two-inertial-reduction}, it suffices to induct on $q$ for a division algebra achieving the minimal unramified index $n=p^q$.
We replace the division algebra accordingly and fix $R=\mathcal O_F$, $P=\mathcal O_D$, a uniformizer $\pi$ of $F$, and $J_D=J(P)$.
Lemma~\ref{lem:review-minimal-order-k1-support} gives ramification index $n$, residue field $\Delta=k^{1/n}$, and period $n$ over the maximal unramified field.

\subsubsection{Two cyclic covers that lower the index}
To run the induction on the minimal unramified index, we construct, after a suitable base extension, a cyclic extension $T/F$ of degree $p$ inside $D$ whose centralizer has smaller minimal unramified index.
There are two types of such extensions, according to the ramification of $T/F$.

\begin{lemma}\label{lem:second-two-cyclic-cover}
Let $k$ be a nonarchimedean local field of characteristic $p$, let $F$ be a complete discretely valued field with first residue field $k$, and let $D/F$ be a central division algebra with $\deg D=\ind D_{F^{\mathrm{nr}}}=n=p^q>1$.
Put $m=n/p$, $P=\mathcal O_D$ and $\Delta=P/J(P)$.
After adjoining $\mu_p$ if necessary and passing to a finite Galois unramified extension, and renaming the resulting fields and orders again $F$, $k$, $D$, $P$, $\Delta$ and $\pi$, there are a cyclic extension $T/F$ of degree $p$, a generator $\sigma$ of $\Gal(T/F)$ and $x\in D^\times$ with the following properties.
\begin{enumerate}
\item $T\subset D$, and conjugation by $x$ induces $\sigma$ on $T$.
\item $B=Z_D(T)$ is a division algebra with $\deg B=\ind B_{T^{\mathrm{nr}}}=m$; it has ramification index $m$, its residue field $\Delta_0$ is the field of $m$-th roots of the residue field of $T$, and $\SK_1(B)=0$.
\item $m[D]=(T/F,\sigma,\Nrd_{B/T,1}(x^p))$, where $\sigma$ is the generator in the chosen cyclic normal form.
\item Either $T/F$ is totally ramified and $x$ is a unit whose residue generates $\Delta$ over $k$ \textup{(}the first type, or homogeneous-unit type\textup{)}, or $T/F$ has ramification index one and purely inseparable residue degree $p$, and $x$ is a uniformizer of $D$ \textup{(}the second type, or homogeneous-uniformizer type\textup{)}.
\item With $C=\mathcal O_B$, one has $P=\bigoplus_{j=0}^{p-1} Cx^j$.
\end{enumerate}
\end{lemma}
\begin{proof}
Adjoining $\mu_p$ is an extension of degree prime to $p$, so it preserves the $p$-primary index, and restriction to it detects the $p$-primary part of $\SK_2$.
It also preserves the minimal unramified index, since the maximal unramified extension after adjoining is $F(\mu_p)F^{\mathrm{nr}}$, of degree prime to $p$ over $F^{\mathrm{nr}}$.
The finite Galois unramified stage is not assumed to have prime degree; we return from it by the inclusion through transfer in Lemma~\ref{lem:second-two-inertial-reduction}.
Modulo classes killed by the maximal unramified extension, the top component $m[D]$ has one of the following presentations using a cyclic extension $T/F$ of degree $p$:
\begin{enumerate}
\item $T/F$ is totally ramified, and the coefficient is a unit whose residue is not a $p$-th power.
\item $T/F$ has ramification index one and purely inseparable residue degree $p$, and the coefficient is a uniformizer.
\end{enumerate}
In equal characteristic, the leading pole order in the Artin--Schreier--Cartier presentation gives these two types; see the proof of \cite[\S4, Lemma 5, pp.~338--339]{KatoLocalI}.
In mixed characteristic, depth zero and positive depths prime to $p$ give the first type, and positive depths divisible by $p$ give the second \cite[\S2, Proposition 2(i)--(iii), pp.~321--322, and \S4, p.~339]{KatoLocalI}; at depth zero, Kato's symbol $\{f,\pi\}$ is used in the form $\{\pi,f^{-1}\}$, which gives the totally ramified Kummer field $T=F(\pi^{1/p})$.
Indeed, at an intermediate depth $s$, write $z^p=1+c$ with $v_F(c)=s<pv_F(p)/(p-1)$.
If $p\nmid s$, then $v(z-1)=s/p$.
If $s=ph$, then $Y=(z-1)/\pi^h$ satisfies $\overline Y^{\,p}=\overline{c/\pi^s}$: the intermediate binomial coefficients are divisible by $p$, and the corresponding terms vanish modulo $\pi$ because $v_F(p)-(p-1)h>0$.
In the second type, $\overline{c/\pi^s}\notin k^p$, so $\overline Y$ is a purely inseparable residue root.
At the terminal depth $pv_F(p)/(p-1)$, we obtain a separable Artin--Schreier equation by subtracting $1$ from a $p$-th root and dividing by $\zeta_p-1$, and hence that layer is killed unramifiedly.

In both types $T^{\mathrm{nr}}=TF^{\mathrm{nr}}$, since the residue extension of $T/F$ is trivial or purely inseparable, so the separable residue extensions come from $k$.
Both types of $T/F$ have no nontrivial unramified subextension, so $[T^{\mathrm{nr}}:F^{\mathrm{nr}}]=p$.
After removing the difference that vanishes over the maximal unramified extension, $m[D]$ is a cyclic class for $T$, hence vanishes over $T^{\mathrm{nr}}$ and over some finite stage $TF'$.
The upper bound is
\[
 \ind D_{T^{\mathrm{nr}}}\le\ind D_{TF'}=\per D_{TF'}\le m,
\]
where $F'/F$ is a finite unramified extension such that $m[D]$ splits over $TF'$, and the equality holds over the two-dimensional local field $TF'$ \cite[Corollary 5.1, Proposition 5.2]{Chipchakov2022}.
The lower bound $\ind D_{T^{\mathrm{nr}}}\ge\per D_{T^{\mathrm{nr}}}\ge m$ follows from $\per D_{F^{\mathrm{nr}}}=n$ and $\cor_{T^{\mathrm{nr}}/F^{\mathrm{nr}}}\res_{T^{\mathrm{nr}}/F^{\mathrm{nr}}}=p$ on $\Br(F^{\mathrm{nr}})$.
Hence the minimal unramified index after restriction to $T$ is $m$.
Since $T^{\mathrm{nr}}=TF^{\mathrm{nr}}$, Lemma~\ref{lem:second-two-inertial-reduction} lets us pass to a further finite unramified stage of $F$ at which $\ind D_T=m$ and the difference from the normal form vanishes.
The index of $D$ remains $n$, and since $\ind D/\ind D_T=n/m=[T:F]$, the field $T$ embeds into $D$.
Then $B=Z_D(T)$ is a division algebra of minimal unramified index $m$, and Lemma~\ref{lem:review-minimal-order-k1-support} applied to $B/T$ shows that $B$ has ramification index $m$ and residue field equal to the full field of $m$-th roots of the residue field of $T$.
We have $\per B=\ind B$, and Corollary~\ref{cor:review-two-local-sk1-zero} gives $\SK_1(B)=0$.

For a generator $\sigma$ of $\Gal(T/F)$ and $x\in D^\times$ implementing $\sigma$ by conjugation,
\[
 m[D]=(T/F,\sigma,\Nrd_{B/T,1}(x^p)).
\]
This identity is obtained in the corner of $D^{\otimes_F m}$ corresponding to $T^{\otimes_F m}\to T$.
The algebra $B^{\otimes_T m}$ is split.
For the standard representation $V$ of $\operatorname{GL}_m$, scalars act trivially on $V^{\otimes m}\otimes\det(V)^{-1}$, so $g\mapsto g^{\otimes m}/\det g$ is a representation of $\operatorname{PGL}_m$ and applies functorially to the cocycle and to the $\sigma$-semilinear action; hence the carry in the semilinear action on that corner is $(x^p)^{\otimes m}/\Nrd_{B/T,1}(x^p)$.
If the semilinear action is $P\sigma$, replacing the corner generator $x'$ by $P^{-1}x'$ changes it to componentwise $\sigma$, and the remaining coefficient is $(P^{-1}x')^p=\Nrd_{B/T,1}(x^p)$.

In the first type, correction by $B^\times$ makes $x$ a unit.
Residues of norms of units in a totally ramified extension are $p$-th powers, and $\overline{\Nrd_{B/T,1}(x^p)}=\bar x^{pm}=\bar x^n$ because the residue field $k^{1/m}$ of $B$ is purely inseparable of degree $m$ over $k$.
Hence $\bar x^n\notin k^p$, and $\bar x$ generates $\Delta$.
In the second type, $\Nrd_{B/T,1}(x^p)/\pi$ is a norm from $T/F$, whose valuation lies in $p\Z$, and $v_F(\Nrd_{B/T,1}(x^p))=v_D(x)$, with the valuation of a uniformizer of $D$ normalized to one; hence $v_D(x)\equiv1\bmod p$.
Since $v_D(B^\times)=p\Z$, we may take $x$ to be a uniformizer.
Put $C=\mathcal O_B$ and $\Delta_0=C/J(C)$.
In both types, $P=\bigoplus_{j=0}^{p-1} Cx^j$.
In the first type, after cancelling the minimum valuation of the coefficients by an element of $B$, the independence of $1,\bar x,\ldots,\bar x^{p-1}$ over $\Delta_0$ gives the decomposition; in the second, the valuations of the terms are distinct modulo $p$.
\end{proof}

We pass to the extension of Lemma~\ref{lem:second-two-cyclic-cover}, rename the resulting fields and orders again $F$, $k$, $D$, $P$, $\Delta$ and $\pi$, and fix $T$, $\sigma$, $x$, $B$, $C$ and $\Delta_0$ as in that lemma.
Induction on the minimal unramified index gives $\SK_2(B)=0$.
Together with the isomorphism $\Nrd_2|_{I_T(B)}\colon I_T(B)\to K_2(T)$ of Proposition~\ref{prop:review-splitting-transfer-norm-iso}, this yields $K_2(B)=I_T(B)$.

\subsubsection{The degree-zero relative trace}
The cyclic extension of Lemma~\ref{lem:second-two-cyclic-cover} makes $P$ a $\Z/p$-graded order.
In this subsubsection, we show that the map from the $p$-completed $\mathrm{TC}$ of the total-degree-zero part of $\mathrm{THH}(P)$ to $\mathrm{TC}(P)^\wedge_p$ is surjective on $\pi_2$, and we then remove the residue contribution.
Let $G=\Z/p$.
For a $G$-graded ring $S$, its cover $\mathcal Q(S)$ has vertices $i\in G$, morphisms $S_{j-i}$, and the strict $G$-action that simultaneously translates the vertices.
We write $K(\mathcal Q(S))$ for the $K$-theory of the category of perfect right $\mathcal Q(S)$-modules.
We write $\mathbb T$ for the circle group.

\begin{lemma}
\label{lem:review-cyclic-zero-tc}
Let $A$ be a finite abelian group with $pA=0$, and let $S=\bigoplus_{g\in A}S_g$ be an $A$-graded discrete ring.
Let
\[
 X=\mathrm{THH}(S)=\bigoplus_{g\in A}X_g
\]
be the total-degree decomposition of the cyclic bar construction.
This decomposition preserves the $\mathbb T$-action, and $X_0$ is a $p$-cyclotomic subobject.
We call $\mathrm{TC}(X_0;p)^\wedge_p$ the degree-zero $\mathrm{TC}$ of $S$; we do not claim that it is a direct summand of $\mathrm{TC}(S)^\wedge_p$.
With this notation there is a fiber sequence
\[
 \mathrm{TC}(X_0;p)^\wedge_p\longrightarrow
 \mathrm{TC}(S)^\wedge_p\longrightarrow
 \left(\bigoplus_{g\ne0}\Sigma(X_g)_{h\mathbb T}\right)^\wedge_p.
\]
If $\pi_1(((X_g)_{h\mathbb T})^\wedge_p)=0$ for all $g\ne0$, the first map is therefore surjective in degree two.

For the graded cover $\mathcal Q_A(S)$ with vertices $A$ and morphisms $S_{j-i}$ from $i$ to $j$, there is a natural equivalence
\[
 \left(\mathrm{TC}(\mathcal Q_A(S);p)_{hA}\right)^\wedge_p
 \simeq\mathrm{TC}(X_0;p)^\wedge_p.
\]
On the left, $A$ acts on $\mathcal Q_A(S)$ by vertex translation.
The same assertions hold for relative fibers of maps of graded rings or covers.
\end{lemma}
\begin{proof}
Total degree in the cyclic bar construction is preserved by face maps given by multiplication, degeneracies given by the unit, and cyclic permutation.
The cyclotomic Frobenius from the $p$-fold Tate diagonal repeats a simplex of total degree $g$ exactly $p$ times, sending it to degree $pg=0$.
Thus $X_0$ is a sub-$p$-cyclotomic spectrum, and the cofiber $Q=X/X_0$, with underlying spectrum $\bigoplus_{g\ne0}X_g$, has zero Frobenius.

Since $Q$ has zero Frobenius, the Nikolaus--Scholze formula for bounded-below $p$-cyclotomic spectra identifies $\mathrm{TC}(Q;p)^\wedge_p$ with the fiber of the canonical map $(Q^{h\mathbb T})^\wedge_p\to(Q^{t\mathbb T})^\wedge_p$.
By the circle norm fiber sequence, we have $\mathrm{TC}(Q;p)^\wedge_p\simeq(\Sigma Q_{h\mathbb T})^\wedge_p$.
The functor $\mathrm{TC}(-;p)$ is exact between stable categories; applying it to $X_0\to X\to Q$ gives the stated sequence.
Since circle homotopy orbits commute with the finite direct sum $\bigoplus_{g\ne0}X_g$, we have $Q_{h\mathbb T}\simeq\bigoplus_{g\ne0}(X_g)_{h\mathbb T}$.
The degree-two assertion follows from the long exact sequence \cite[Proposition II.1.9, Lemma II.4.2]{NikolausScholzeTC}.

In each simplicial degree, the lifts of a total-degree-zero simplex in the cyclic bar construction form a free $A$-set indexed by the starting vertex.
Taking homotopy orbits under translation of the starting vertex gives $\mathrm{THH}(\mathcal Q_A(S))_{hA}\simeq X_0$ via the map that forgets the vertices.
For every prime $q$, the $q$-fold Tate diagonal repeats a cycle $q$ times from the same starting vertex and commutes with translation of vertices, so the $A$-action and this map are cyclotomic \cite[Corollary II.1.7, \S III.2]{NikolausScholzeTC}.
Since the forgetful functor from cyclotomic spectra preserves colimits, the computation of free $A$-orbits in each cyclic degree gives an equivalence of cyclotomic spectra.
The absolute $\mathrm{THH}$ is connective, and $\mathrm{TC}/p$ preserves all colimits of connective cyclotomic spectra \cite[Theorem 2.7]{ClausenMathewMorrow}, so it may be applied to the $A$-bar construction.
Induction using the coefficient exact sequence gives the same assertion with Moore coefficients $p^a$, and passing to the inverse limit over $a$ yields the displayed equivalence in the $p$-complete category.
The relative version follows from the two absolute versions, because finite fibers commute with homotopy orbits and $\mathrm{TC}$ is exact.
\end{proof}

\begin{lemma}\label{lem:review-finite-pbasis-differentials}
Let $R$ be a complete discrete valuation ring with uniformizer $\pi$ and residue field $k$ of characteristic $p$.
Let $\{\bar u_i\}$ be a finite $p$-basis of $k$, and choose lifts $u_i\in R$.
In mixed characteristic, $\Omega^1_{R/\Z}\otimes_R R/p^aR$ is generated by $d\pi$ and the $du_i$ for every $a\ge1$.
Moreover, for $f\in\pi^tR$,
\[
 df\in\pi^{t-1}R\,d\pi+
       \sum_i\pi^tR\,du_i
 \quad\pmod {p^a\Omega^1_{R/\Z}}.
\]
In equal characteristic, choose a coefficient field and take the $u_i$ in it.
Then the $u_i$ together with $\pi$ form a $p$-basis of the ring $R$, and $\Omega^1_{R/\F_p}$ is free on $d\pi$ and the $du_i$.
\end{lemma}
\begin{proof}
In mixed characteristic, put $e=v_R(p)$ and choose $s\ge a$ and $M\ge ea+1$.
By lifting the residue field's $p^s$-basis expansion one digit at a time, we express every $r\in R$ as
\[
 r=\sum_{t=0}^{M-1}\pi^t
      \sum_{0\le I_i<p^s}b_{t,I}^{p^s}u^I+\pi^Mc
 \qquad(b_{t,I},c\in R).
\]
For any derivation $d$ of $R$ into a module killed by $p^a$, we have $d(b^{p^s})=0$ and $d(\pi^Mc)=0$.
Thus $dr$ lies in the submodule generated by the stated finite collection of differentials.
For $f=\pi^tr$, the Leibniz rule gives $df=t\pi^{t-1}r\,d\pi+\pi^t\,dr$, which proves the depth assertion.
This argument applies to universal differentials, so no finite-type hypothesis on $R$ is needed.

In equal characteristic, write $R=k[[\pi]]$.
Using the $p$-basis expansion of coefficients and the decomposition of power-series exponents according to their remainders modulo $p$, we obtain
\[
 R=\bigoplus_{0\le I_i<p,\ 0\le j<p}
                  R^p u^I\pi^j.
\]
This is precisely the definition of a ring $p$-basis, proving the assertion about differentials.
\end{proof}

\begin{lemma}
\label{lem:review-finite-free-pro-kernel}
Let $R$ be a discrete valuation ring with uniformizer $\varpi$, and let $L,L'$ be finite free $R$-modules.
If $f\colon L\to L'$ is injective, there is $c\ge0$ such that, for all $j\ge0$ and all $s\ge j+c$,
\[
 \ker(f\bmod\varpi^s)\longrightarrow L/\varpi^jL
\]
is zero.
\end{lemma}
\begin{proof}
We clear denominators in a left inverse over the fraction field and obtain $g\colon L'\to L$ and $c\ge0$ with $gf=\varpi^c\mathrm{id}_L$.
If $a\in L$ satisfies $f(a)\in\varpi^sL'$, then $\varpi^ca\in\varpi^sL$ and hence $a\in\varpi^{s-c}L\subseteq\varpi^jL$.
\end{proof}

\begin{lemma}
\label{lem:review-nonzero-hc-one}
For each nonzero degree $g$ of the $\Z/p$-graded maximal order $P$ constructed above, one has, in mixed characteristic,
\[
 \{HC_1(P/p^aP\,/\,\Z/p^a)_g\}_{a\ge1}=0
 \quad\text{in the pro category}
\]
and, in equal characteristic, $HC_1(P/\F_p)_g=0$.
\end{lemma}
\begin{proof}
We normalize the valuation on the center by $v_D(\pi)=n$.
The ring $\gr D$ is commutative.
Indeed, since $\Delta/k$ is purely inseparable, conjugation by a homogeneous element induces the identity $k$-automorphism of $\Delta$.
The quotient of value groups is cyclic, and every homogeneous component is the product of $\Delta$ and a power of one uniformizer; these commute.
We calculate cyclic one-forms in $\Omega^1(P)_\natural=\Omega^1(P)/[P,\Omega^1(P)]$ and then subtract actual exact differentials.
The boundary is $b(y\,dz)=[y,z]$.

We begin with the homogeneous-unit type.
Let $\eta$ be a degree-zero uniformizer and $x$ a degree-one unit.
Residue and value-group bases give the $R$-basis $\eta^i x^j$ ($0\le i,j<n$) of $P$.
We choose the lift $u$ of the central $p$-basis element as the constant term of the minimal polynomial of $x$.
Then
\[
 \varphi_\eta(\eta)=\eta^n+\pi\sum_{i<n}b_i\eta^i=0,
 \quad b_0\in R^\times,\qquad
 \varphi_x(x)=x^n+\sum_{0<j<n,\,p\mid j}c_jx^j-u=0,
 \quad c_j\in\pi R.
\]
Only exponents divisible by $p$ occur in $\varphi_x$, here and in the homogeneous-uniformizer type below: if $f$ is the minimal polynomial of $x^p$ over $F$, then $f(X^p)$ kills $x$ and has degree $[F(x):F]$, so it is the minimal polynomial of $x$.
Lemma~\ref{lem:review-finite-pbasis-differentials} expresses every coefficient differential on a finite truncation in terms of $d\pi,du$.

For $y$ of nonzero degree, we apply $y\,d\varphi_\eta(\eta)=0$.
Let $\psi(y)$ be the coefficient of $d\eta$ after cyclic permutation.
Then $[\psi(y),\eta]=[y,\varphi_\eta(\eta)]=0$.
The field $F(\eta)$ is a self-centralizing maximal subfield in degree zero, so $\psi(y)=0$.
The relation therefore involves only $d\pi$ and $du$.
The coefficient operator of $d\pi$ is multiplication by the unit $\sum_i b_i\eta^i$ plus $\pi$ times an integral operator, hence is invertible.
The coefficient of $du$ is divisible by $\pi$.
A $d\pi$ coefficient of depth $r$ can therefore be replaced by a $du$ coefficient of depth at least $r+n$.

Next we use $y\,d\varphi_x(x)=0$ and eliminate $d\pi$ as above.
The coefficient operator of $du$ is $-1$ plus $\pi$ times an integral operator, hence is invertible.
For the coefficient of $dx\,x^{-1}$, each cyclic sum on the associated graded ring is multiplication by the corresponding exponent $j$.
Every such $j$ is divisible by $p$, so this coefficient is at least one step deeper than the input.
Thus a $du$ coefficient of depth $r$ can be replaced by a $dx\,x^{-1}$ coefficient of depth at least $r+1$; for $d\pi$, the gain is at least $n+1$.

We take the exact differential of $h=f\eta^i x^j$, where $f\in\pi^tR$ and $j\equiv g\pmod p$.
After scalar differentials are eliminated by the two relations above, the coefficient of $dx\,x^{-1}$ is $j h$ modulo depth $v_D(h)+1$.
Indeed, the $d\pi$-part of $df$ loses at most $n$ in depth, but its replacement gains $n+1$.
The $du$-part loses no depth and gains one step.
The cyclic sum in $d(x^j)$ is multiplication by $j$ in the commutative associated graded ring.
Since $j\equiv g\ne0$ is a unit, subtracting such actual exact differentials depth by depth eliminates the $dx\,x^{-1}$ coefficient.
Modulo $p^a$ this stops after finitely many steps, leaving a form $z\,d\eta$.

We apply the same procedure to the homogeneous-uniformizer type.
Let $w$ be a degree-zero unit and $x$ a degree-one uniformizer.
The basis is now $w^i x^j$ ($0\le i,j<n$), with relations
\[
 \varphi_w(w)=w^n+\sum_{i<n}c_iw^i-u=0,
 \quad c_i\in\pi R,\qquad
 \varphi_x(x)=x^n+\pi\sum_{i<m}b_i x^{pi}=0,
 \quad b_0\in R^\times.
\]
For an input $y$ of nonzero degree, the $dw$ coefficient of $y\,d\varphi_w(w)$ commutes with $w$.
Since $F(w)$ is a self-centralizing maximal subfield in degree zero, this coefficient vanishes integrally as well.
On the other hand, the $d\pi$ coefficient of $y\,d\varphi_x(x)$ is multiplication by a unit plus $\pi$ times an integral operator.
The $dx$ coefficient of $y\,d\varphi_x(x)$ increases the depth of $y$ by at least $n$.
In the highest term, $x^{n-1}$ contributes $n-1$, and the cyclic sum of length $n$, divisible by $p$, contributes one further step.
The lower terms satisfy the same bound because of their explicit factor $\pi$ and the cyclic sums for exponents $pi$.

These two relations therefore eliminate $d\pi,du$ simultaneously, moving a scalar-differential coefficient of depth $r$ to a $dx$ coefficient of depth at least $r+n$.
The leading $dx$ coefficient of the exact differential of $h=fw^i x^j$ is $j h x^{-1}$.
The $d\pi$-part of $df$ loses at most $n$, so after elimination it is at least one step deeper than the leading depth $v_D(h)-1$.
Again $j\equiv g\ne0$ allows successive elimination of $dx$ by exact differentials, leaving $z\,dw$.

Thus, for both types, every $HC_1$ class on a finite truncation is represented respectively by
\[
 z\,d\eta,\quad [z,\eta]=0\pmod {p^aP};
 \qquad z\,dw,\quad [z,w]=0\pmod {p^aP}.
\]
The coefficient $z$ belongs to the finite free module $P_g/p^aP_g$.
The maps $[\,\cdot\,,\eta]$ and $[\,\cdot\,,w]$ are injective on $P_g$, because the self-centralizing subfields lie in degree zero.
By Lemma~\ref{lem:review-finite-free-pro-kernel}, for any fixed truncation, these kernels at a sufficiently deep stage have zero image in that truncation.
The inverse system of $HC_1$ groups is therefore pro-zero.

In equal characteristic, $\Omega^1_{R/\F_p}$ is finite free and the cyclic one-forms form a finitely generated, hence complete, $R$-module.
The primitives chosen successively above have depths tending to infinity, so they converge.
Continuity of differentiation lets us subtract their sum as one exact differential.
Integral injectivity of the remaining boundary makes $HC_1$ itself zero.
\end{proof}

This lemma supplies the vanishing needed for comparison with absolute degree-zero $\mathrm{TC}$.

\begin{lemma}\label{lem:second-two-degree-zero-surjective}
For the $\Z/p$-graded maximal order $P$ of Lemma~\ref{lem:second-two-cyclic-cover}, with $X=\mathrm{THH}(P)$, the map
\[
 \mathrm{TC}(X_0;p)^\wedge_p\longrightarrow\mathrm{TC}(P)^\wedge_p
\]
of Lemma~\ref{lem:review-cyclic-zero-tc} is surjective on $\pi_2$.
\end{lemma}
\begin{proof}
For each $g\ne0$, we compute the completion of the circle orbits $Z_g=(\mathrm{THH}(P)_g)_{h\mathbb T}$.
In mixed characteristic, $P$ is flat over $\Z$, and low-degree linearization and base change give
\[
 \pi_1(Z_g/p^a)\simeq
 HC_1(P/p^aP\,/\,\Z/p^a)_g.
\]
Here we use the degree-one comparison of circle orbits for $\mathrm{THH}(P)\to HH(P/\Z)$.
In the term of simplicial degree $s$ of the cyclic bar construction, the map from the topological tensor product to the tensor product over $\Z$ is an isomorphism on $\pi_0$ and $\pi_1$: since $P$ is $\Z$-flat, the classical terms are concentrated in degree zero, and the topological ones have $\pi_0=P^{\otimes(s+1)}$ and $\pi_1=0$.
This low-degree comparison is preserved by Moore reduction, realization and circle orbits, and the classical reduction modulo $p^a$ is the cyclic complex of $P/p^aP$.
Lemma~\ref{lem:review-nonzero-hc-one} therefore makes the inverse system on the left pro-zero.
In equal characteristic, the fiber of $\mathrm{THH}(\F_p)\to H\F_p$ lies in degrees at least two, so $\pi_1Z_g\simeq HC_1(P/\F_p)_g=0$ by Lemma~\ref{lem:review-nonzero-hc-one}.
The Moore sequence identifies the transitions on $\pi_1(Z_g/p^a)=\pi_0Z_g[p^a]$ with multiplication by $p$.
Since $p\pi_0Z_g=0$, these transitions are zero, so the inverse system is pro-zero in this case as well.

In either case, the inverse system $\{\pi_1Z_g/p^a\}_a$ has surjective transitions and embeds into the preceding pro-zero system, so every term is zero.
Thus $\pi_1Z_g$ is $p$-divisible, and the maps $\pi_1Z_g[p^{a+1}]\xrightarrow{p}\pi_1Z_g[p^a]$ are surjective.
Both outer terms of the Moore exact sequence
\[
 0\to\pi_2Z_g/p^a\to\pi_2(Z_g/p^a)\to \pi_1Z_g[p^a]\to0
\]
have surjective transitions, so the middle term satisfies Mittag--Leffler as well.
The Milnor exact sequence and $\varprojlim_a\pi_1(Z_g/p^a)=0$ give $\pi_1((Z_g)^\wedge_p)=0$.
By Lemma~\ref{lem:review-cyclic-zero-tc}, the $\mathrm{TC}$ of the cofiber of the degree-zero subobject is $(\bigoplus_{g\ne0}\Sigma Z_g)^\wedge_p$, whose degree two vanishes.
This gives surjectivity in degree two of the map from absolute degree-zero $\mathrm{TC}$ to $\mathrm{TC}(P)^\wedge_p$.
\end{proof}

The lift provided by Lemma~\ref{lem:second-two-degree-zero-surjective} need not vanish in the residue orbits $(\mathrm{TC}(\bar H)^\wedge_p)_{hG}$, where $\bar H$ is the semisimple residue of the cover $H=\mathcal Q(P)$; removing this component makes the remaining class relative, so that Lemma~\ref{lem:app-orbit-stable-image} applies.
In the homogeneous-uniformizer type, $\mathcal Q(P)$ is the Iwahori order with $C$ on and above the diagonal and $J(C)$ below it.
The group $G$ regularly permutes the $p$ copies of $\Delta$ in the semisimple residue of $\mathcal Q(P)$, so the residue orbit is $\mathrm{TC}(\Delta)^\wedge_p$.
The groups $\pi_2\mathrm{TC}(\Delta)^\wedge_p$ and $\pi_3\mathrm{TC}(\Delta)^\wedge_p$ vanish because $[\Delta:\Delta^p]=p$ \cite[Theorem 6.5]{ClausenMathewMorrow}, \cite[Theorem 8.1]{GeisserLevine}.
In the homogeneous-unit type, put $E=F(x)$.
Then $E/F$ is a maximal subfield of ramification index one and residue field $\Delta$.
The ring $\mathcal O_E=\bigoplus_{i<p}\mathcal O_{F(x^p)} x^i$ defines a graded subcover.
In equal characteristic, choose a coefficient field $\Delta\subset\mathcal O_E$ containing $x$, where $\bar x$ is a $p$-basis of $\Delta$; since $E^p\subset F(x^p)$, the field $\Delta_0=\Delta^p$ lies in the degree-zero part $\mathcal O_{F(x^p)}$, so the residue cover has a graded section.
Once the cover is forgotten, the subtracted residue term is zero because $\pi_2\mathrm{TC}(\Delta)^\wedge_p=0$.

We now treat mixed characteristic.
Put $S=\mathcal O_{F(x^p)}$, $l=\Delta_0$, and $v=x^p\in S^\times$.
The reduction of the commutative subcover is the same graded cover of $\Delta=\bigoplus_{i<p}\Delta_0\bar x^i$ as the reduction of $\mathcal Q(P)$.
Multiplication by $x$ identifies successive vertices with the corner.
On $\operatorname{Perf}(S)$ the transported action of $G=C_p$ is the identity for each group element.
For $0\le i,j<p$, writing $[i+j]$ for the representative in $\{0,\ldots,p-1\}$, its composition isomorphism is multiplication by
\[
 x^{i+j-[i+j]}=v^{\lfloor(i+j)/p\rfloor}.
\]
The action after reduction is given by the same formula with $\bar v\notin l^p$.
We show that reduction on the $p$-completed homotopy orbits of $\mathrm{TC}(S)$ is surjective on $\pi_2$.
For this we compare the trace with the determinant.
For a finite separable $l'/l$, let $S_{l'}$ be the corresponding finite \'etale local $S$-algebra.
On the small \'etale site of $l$, let $\mathcal D_S$ and $\mathcal D_l$ be the Picard spectra of graded lines over the algebras $S_{l'}$ and over the fields $l'$, respectively, with the symmetry on ranks $r,r'$ given by $(-1)^{rr'}$.
For $\Lambda\in\{S,l\}$, put $\mathcal D(\Lambda)=R\Gamma(l,\mathcal D_\Lambda)^\wedge_p$.
The group $G$ acts on these spectra through the action above, with the composition isomorphisms $v_\Lambda^{\lfloor(i+j)/p\rfloor}$, which act on determinants of rank $n$ by $v_\Lambda^{n\lfloor(i+j)/p\rfloor}$, where $v_S=v$ and $v_l=\bar v$.

\begin{lemma}\label{lem:second-two-determinant-comparison}
In this setup, the determinant and the cyclotomic trace give a map of reduction arrows of $G$-spectra
\[
 \bigl(\mathrm{TC}(S)^\wedge_p\to\mathrm{TC}(l)^\wedge_p\bigr)
 \longrightarrow\bigl(\mathcal D(S)\to\mathcal D(l)\bigr).
\]
The map $\mathrm{TC}(l)^\wedge_p\to\mathcal D(l)$ is an equivalence, and the map $\bigl((\mathrm{TC}(S)^\wedge_p)_{hG}\bigr)^\wedge_p\to\bigl(\mathcal D(S)_{hG}\bigr)^\wedge_p$ is surjective on $\pi_2$.
\end{lemma}
\begin{proof}
We first take the signed determinant of perfect complexes and then take Moore coefficients.
This determinant identifies the zero and first homotopy sheaves, and a scalar automorphism $v$ on a complex of rank $n$ induces multiplication by $v^n$ on its determinant \cite[Proposition 5.3 and \S12]{BhattScholzeWitt}.
The finite \'etale base-change functors preserve the scalar composition isomorphisms.
The cyclotomic trace is natural for these functors and for their natural isomorphisms \cite[Theorem 7.4]{BlumbergGepnerTabuada}, and so is the determinant \cite[Proposition 5.3]{BhattScholzeWitt}.
Over the closed \'etale site, the cyclotomic trace identifies the sheafified Postnikov towers of $K/p^a$ and $\mathrm{TC}/p^a$ by strict Henselian comparison.
For each $\Lambda\in\{S,l\}$, the global sections of the latter tower compute $\mathrm{TC}(\Lambda)/p^a$ by proper base change for the identity of $\operatorname{Spec}\Lambda$, since $\Lambda$ is Henselian along $(p)$ \cite[Theorems 6.1 and 6.3]{ClausenMathewMorrow}.
Inverting the trace equivalence, composing with the determinant, and taking the inverse limit gives the map of reduction arrows of $G$-spectra in the statement.
For $l$, the determinant is an equivalence with finite Moore coefficients on the geometric stalk.
Indeed, a separable extension of $l$ has $p$-rank one, so logarithmic differentials in degree at least two vanish.
The norm-residue isomorphism and the comparison with Quillen $K$-theory then give $K_i(-;\Z/p^a)=0$ there for $i\ge2$, and the units have no $p$-torsion \cite[Corollary (2.8)]{BlochKato1986}, \cite[Theorem 8.1]{GeisserLevine}.
It follows by descent that $\mathrm{TC}(l)^\wedge_p\to\mathcal D(l)$ is an equivalence.

Next we bound the comparison fiber.
Write $U$ for the sheaf of principal units $1+\mathfrak m_{S_{l'}}$.
Its higher cohomology on $l$ vanishes.
At a finite Galois stage the successive quotients of the principal-unit filtration are additive residue modules, whose higher cohomology vanishes by the normal basis theorem.
The invariants of the finite congruence quotients form a surjective tower, so the vanishing passes to the inverse limit there, and continuous cochains on the geometric stalk are the filtered union of the finite-stage cochains.
Also $U(S)$ has bounded $p$-power torsion, since a primitive $p^a$-th root of unity satisfies
$v_S(\zeta-1)=v_S(p)/(p^{a-1}(p-1))$.
The rank map is unchanged by reduction, and the scalar coherences act trivially on objects of rank zero.
Thus the fiber of $\mathcal D(S)\to\mathcal D(l)$ is $\Sigma H\widehat{U(S)}_p$ with trivial $G$-action.
The relative group $K_1(S,l)$ is $U(S)$ because $K_2(S)\to K_2(l)$ is surjective by lifting Matsumoto symbols, and $K(S,l)$ is $1$-connective.
The determinant induces this identification on $\pi_1$ of the relative spectra.
The Henselian comparison gives $K(S,l)^\wedge_p\simeq \mathrm{TC}(S,l)^\wedge_p$ \cite[Theorem 4.36]{ClausenMathewMorrow}.
The relative comparison fiber is therefore $2$-connective, also after completion since the bottom Moore homotopy groups form a surjective tower.
The absolute comparison fiber over $S$ has the same bound, since the comparison over $l$ is an equivalence.
Homotopy orbits and subsequent $p$-completion preserve this bound, so the comparison over $S$ is surjective on $\pi_2$ of the completed orbits.
For a bounded-below $p$-complete spectrum, the homotopy orbits are already $p$-complete, since the $G$-bar construction has $p$-complete finite skeleta and the cofibers of the skeleta in its realization have increasing connectivity.
\end{proof}

We compute the reduction map on these determinant orbits.

\begin{lemma}\label{lem:second-two-determinant-orbits}
In this setup, assume that $l$ is a nonarchimedean local field.
Then the reduction map
\[
 \pi_2\bigl((\mathcal D(S)_{hG})^\wedge_p\bigr)\longrightarrow\pi_2\bigl((\mathcal D(l)_{hG})^\wedge_p\bigr)
\]
is surjective.
\end{lemma}
\begin{proof}
Put
\[
 \mathcal A=R\Gamma(l,H\underline{\Z})^\wedge_p,\qquad
 \mathcal M_\Lambda=R\Gamma(l,H\mathbf G_{m,\Lambda})^\wedge_p\quad(\Lambda\in\{S,l\}).
\]
The rank and unit terms in $\Sigma\mathcal M_\Lambda\to\mathcal D(\Lambda)\to\mathcal A$ have trivial action.
Write $\kappa_\Lambda$ for the connecting map with constant action and $\iota_{v_\Lambda}\colon\mathcal A\to\mathcal M_\Lambda$ for the map induced by $n\mapsto v_\Lambda^n$, where $v_S=v$ and $v_l=\bar v$.
The connecting map for the weak action is
\[
 \kappa_\Lambda+(\Sigma^2\iota_{v_\Lambda})\circ e_G(\mathcal A),
 \qquad e_G(\mathcal A)\colon\mathcal A\longrightarrow\Sigma^2\mathcal A,
\]
where $e_G(\mathcal A)$ is the tensor product of the standard double extension $H\Z_p\to\Sigma^2H\Z_p$, given by the class $e_G\in H^2(G,\Z)$ represented by $(g^i,g^j)\mapsto\lfloor(i+j)/p\rfloor$, with the rank module $\mathcal A$ with trivial action.
Indeed, before taking sections, the space of extensions with identified endpoints is the mapping space from $H\underline{\Z}$ to $\Sigma^2H\mathbf G_m$.
It is $2$-truncated, and the two actions agree on vertices and edges.
Their difference on each $2$-simplex is the scalar determinant $v_\Lambda^{n\lfloor(i+j)/p\rfloor}$.
The same calculation in the category of reduction arrows gives the constant-action comparison of the sign terms and the compatible carry terms.
For $p=2$, the residue sign is trivial because $-1=1$ in $l$, so $\kappa_l=0$.
For odd $p$, tensor square of the line gives a map to the unsigned graded-line spectrum that is an equivalence after $p$-completion, and again $\kappa_l=0$.
We retain $\kappa_S$ as a map of spectra.

We now determine the degree-two orbits.
The two nonzero homotopy groups of $\mathcal A$ are $\Z_p$ in degree zero and $X_l=H^1_{\mathrm{cont}}(l,\Z_p)$ in degree $-1$.
Those of $\mathcal M_l$ are $\widehat{l^\times}_p$ and $T_p\operatorname{Br}(l)=\Z_p$ in the same degrees.
These descriptions follow from $\operatorname{cd}_p(l)=1$, Hilbert~90 and the injectivity of the $p^a$-power map on the \'etale unit sheaf of $l$.
Reduction on $\mathcal M_\Lambda$ is surjective on $\pi_0$ and an isomorphism on $\pi_{-1}$ because its fiber is $H\widehat{U(S)}_p$.
The map $\iota_{\bar v}$ in degree zero sends $1$ to a primitive element of $\widehat{l^\times}_p$ because $\bar v\notin l^p$.
In degree $-1$ it is cup product with $\bar v$.
The local invariant pairing makes this map surjective modulo $p$, and $\operatorname{cd}_p(l)=1$ lifts degree-one characters to $\Z_p$-characters, so the map to $T_p\operatorname{Br}(l)$ is surjective \cite[\S5]{Kurihara2000}.

Put $X_l^0=\ker(X_l\xrightarrow{\cup\bar v}T_p\operatorname{Br}(l))$ and $U^0=\widehat{l^\times}_p/\Z_p\bar v$.
The modules $\widehat{l^\times}_p$ and $X_l$ are $p$-complete and $p$-torsion-free.
For the units this follows from the injectivity of multiplication by $p$ on $l^\times$, and for $X_l$ from the surjective reduction maps on $H^1(l,\Z/p^a)$ and the coefficient exact sequence.
A $p$-complete, $p$-torsion-free $\Z_p$-module is the completion of a free module obtained by lifting a basis modulo $p$.
Maps from the resulting Eilenberg--Mac Lane spectrum in derived $p$-complete $H\Z_p$-modules can thus be computed on those generators before completion.
We can therefore split the two homotopy degrees of $\mathcal A$ and $\mathcal M_l$ and split off the primitive summands in the two degreewise maps of $\iota_{\bar v}$.
Since $e_G(\mathcal A)$ is natural for the chosen decomposition of complete $H\Z_p$-modules, on the residue side, where $\kappa_l=0$, the carry survives only on the corresponding two summands.
Each of the two matched rank and unit summands of the carry is the extension represented by
\[
 0\longrightarrow\Z_p\xrightarrow{\ 1+g+\cdots+g^{p-1}\ }\Z_p[G]
 \xrightarrow{\ g-1\ }\Z_p[G]\longrightarrow\Z_p\longrightarrow0.
\]
Consequently its middle complex is free over $\Z_p[G]$ in degrees $1,0$, and its derived coinvariants are $\Z_p\xrightarrow{0}\Z_p$ in these degrees.
The complex for the degree $-1$ pairing is its shift by $-1$, with derived coinvariants in degrees $0,-1$.
Neither contributes in degree two.
The remaining summands have trivial action, so the degree-two determinant orbits are generated by $H_1(G,U^0)$ and $H_3(G,X_l^0)$.
The unit contribution lifts from $\Sigma\mathcal M_S$ by the surjectivity of unit reduction.

For the character contribution, we restrict the reduction arrow of extensions to the rank summand $\Sigma^{-1}HX_l^0$ of $\mathcal A$.
The restriction of $\iota_v$ is null in derived $p$-complete $H\Z_p$-modules because it is zero on $\pi_{-1}$ and the source is a shifted completed free module.
We choose a nullhomotopy over $S$ and reduce it to obtain the compatible nullhomotopy over $l$.
The carry then disappears on this reduction arrow, leaving the constant-action comparison, including the sign $\kappa_S$ when $p=2$.
Its relative boundary from the character summand lies in the orbits of the constant map
\[
 \Sigma^{-1}HX_l^0\longrightarrow\Sigma^2H\widehat{U(S)}_p .
\]
The contribution from $H_3(G,X_l^0)$ belongs to the reduced summand under $\Sigma^\infty BG_+\simeq\mathbb S\vee\Sigma^\infty BG$.
The constant map preserves this splitting, and $\pi_2$ of the reduced target is $\widetilde H_0(BG,\widehat{U(S)}_p)=0$.
Thus the character contribution lifts as well.
Changing its chosen splitting over $l$ changes its image by unit classes, which already lift.
\end{proof}

By Lemmas~\ref{lem:second-two-determinant-comparison} and~\ref{lem:second-two-determinant-orbits}, the comparison with $\mathrm{TC}$ now lifts every degree-two residue-orbit class from the commutative subcover.
We subtract such a lift to remove the residue image of the absolute degree-zero class.
Once the cover is forgotten, the subtracted term lies in $\pi_2\mathrm{TC}(\mathcal O_E)^\wedge_p$; since $\pi_2\mathrm{TC}(\Delta)^\wedge_p=0$, it lifts to a relative class.
By Corollary~\ref{cor:app-commutative-cut-representative}, the image of this relative class in the prescribed truncation is represented by an integral element of $K_2(\mathcal O_E,\Delta)$ whose generic-point image lies in the image of $\theta_E$.

\subsubsection{Integral transfer from orders and vanishing}
We now complete the proof of Theorem~\ref{thm:second-two-wild}.
Under the inductive hypothesis $K_2(B)=I_T(B)$, we will show in Lemma~\ref{lem:second-two-orbit-image} below that the image of the degree-two homotopy-orbit classes of the relative cover lies in $I_F(D)$.
After representing the image of the relative $\mathrm{TC}$ orbit class at a finite truncation by the trace of an integral class and applying the congruence separation, the injectivity of $\Nrd_2$ on $I_F(D)$ closes the induction.

We consider the relative spectrum $Y=K(H,\bar H)$ of the cover $H=\mathcal Q(P)$, where $\bar H$ is the residue cover used above.
The relative Morita comparison with the corner makes $Y$ $1$-connective, with $\pi_1Y=K_1(C,\Delta_0)$ and action $\tau=\operatorname{Ad}(x)$ of $G$.
In the homogeneous-uniformizer type, the relative $K_1$ classes of $a\in1+J(C)$ placed at different corners agree: for $i<j$ put $b=e_{ij}$ and $d=(a-1)e_{ji}\in J_H=\ker(H\to\bar H)$ and use $[1+bd]=[1+db]$; hence after vertex translation the remaining action is conjugation by the change of coordinates.
Since $\deg B=\ind B_{T^{\mathrm{nr}}}=m$, lifting Matsumoto symbols to a commutative order as in the proof of Lemma~\ref{lem:review-geometric-moore-integral-lift} makes $K_2(C)\to K_2(\Delta_0)$ surjective, so $\pi_1Y=M=\ker[K_1(C)\to\Delta_0^\times]$.
For the corner comparison in the homogeneous-uniformizer type, use the two-term derived Iwahori quotient of Lemma~\ref{lem:app-iwahori-derived-cut}.
Transitions on the corner cofibers of finite truncations vanish in one step.
Hence there is an exact sequence
\[
 (\pi_2Y)_G\longrightarrow \pi_2(Y_{hG})
    \longrightarrow H_1(G,M)\longrightarrow 0.
\]
The generic-point transfer of the left-hand image lies in $I_F(D)$ because $K_2(B)=I_T(B)$.
The right-hand edge can also be realized by integral subcovers, as follows.

\begin{lemma}
\label{lem:review-generic-source-edge-generation}
Let $Y\to Z$ be an equivariant map of spectra for a finite group $G$, with $Y$ $1$-connective and $Z$ connective.
Let $S_\lambda\to Z$ be maps of connective $G$-spectra, let $q\colon\pi_2(Z_{hG})\to \Lambda$ be a homomorphism, and let $I\subset \Lambda$ be a subgroup.
Assume that $q$ sends the image of $(\pi_2Z)_G$ and the entire image of every $\mathcal F_1\pi_2((S_\lambda)_{hG})$ into $I$.
If the image in $H_1(G,\pi_1Z)$ of the edge of $u\in\pi_2(Y_{hG})$ belongs to $\sum_\lambda\im[H_1(G,\pi_1S_\lambda)\to H_1(G,\pi_1Z)]$, then $q$ sends the image of $u$ in $\pi_2(Z_{hG})$ into $I$.
Here $\mathcal F_a$ denotes the filtration of the homotopy-orbit homology spectral sequence.
\end{lemma}
\begin{proof}
For a connective spectrum $X$, the first quotient in degree two is
\[
 \mathcal F_1/\mathcal F_0=E^\infty_{1,1}
 =H_1(G,\pi_1X)/\im
   [d_2\colon H_3(G,\pi_0X)\to H_1(G,\pi_1X)],
\]
and $\mathcal F_0\pi_2(X_{hG})=\im[(\pi_2X)_G\to\pi_2(X_{hG})]$.
Since $Y$ is $1$-connective, $u\in\mathcal F_1\pi_2(Y_{hG})$.
We express its edge image as a finite sum of classes from the specified source $H_1$ groups.
We lift each class in $E^\infty_{1,1}$ to the source $\mathcal F_1\pi_2$.
After their images in $Z$ are subtracted, naturality shows that the remainder lies in $\mathcal F_0\pi_2(Z_{hG})$.
The hypotheses on $q$ now apply.
\end{proof}

As above, we write $R=\mathcal O_F$, $P=\mathcal O_D$, $J_D=J(P)$, $C=\mathcal O_B$, and $J_C=J(C)$.
For $b\in C^\times$, write
\[
 w_b=bx,\qquad E_b=F(w_b).
\]
In the homogeneous-unit type, $\bar b\bar x\notin\Delta_0$, so $\bar b\bar x$ generates $\Delta/k$ and $E_b/F$ is a maximal subfield of residue degree $n$ and ramification index one.
In the homogeneous-uniformizer type, $v_D(w_b)=1$, so $E_b/F$ is a totally ramified maximal subfield of degree $n$.
Since $w_b$ acts nontrivially by conjugation on $T$, we have $w_b\notin B$, but $w_b^p\in B$.
By the Kummer presentation $X^p-w_b^p$ (the purely inseparable presentation in equal characteristic), we have $[E_b:F(w_b^p)]=p$ and $[F(w_b^p):F]=m$.
Moreover, $T\cap E_b=F$, so $T(w_b^p)/T$ is a maximal subfield of $B$.
A residue basis or an Eisenstein basis, respectively, gives
\[
 \mathcal O_{E_b}=\bigoplus_{j=0}^{p-1}\mathcal O_{F(w_b^p)}w_b^j.
\]
This is the commutative graded subcover used below.
For a Brauer class $\beta$ over a two-dimensional local field $L$, put $\chi_\beta(z)=\inv_L(\beta\cup z)$.

\begin{lemma}
\label{lem:review-homogeneous-unit-norm-generation}
With the notation above, let $U_D=1+J_D$, $U_C=1+J_C$, $U_F=1+\mathfrak m_F$, and $V=\Nrd_{B/T}(U_C)$.
For any valued field $L$, write $U(L)=1+\mathfrak m_L$ and $U_T=U(T)$.
In the homogeneous-unit type, the ratios
\[
 \frac{N_{E_b/F}f(bx)}{N_{F(x)/F}f(x)}
\]
for $b\in1+J_C$ and $f(X)=1+hX$ ($h\in R$), together with $N_{T/F}V$, generate $\ker\chi_{[D]}\cap U_F$.
In the homogeneous-uniformizer type, the groups $N_{E_b/F}U(E_b)$ ($b\in C^\times$), together with $N_{T/F}V$, generate the same group.
One can choose $u\in U_F$ such that these generators, together with $u^m=N_{F(x^p)/F}u$, generate $V^G=\ker\chi_{p[D]}\cap U_F$.
\end{lemma}
\begin{proof}
First approximate principal units of $D$ one valuation layer at a time.
In the homogeneous-unit type, let $s$ be a uniformizer of $C$, take $c\in C$, and put $b=1+cs^r$.
The leading coefficient of $f(bx)/f(x)$ at depth $r$ is
\[
 \bar c\,\frac{\bar x f'(\bar x)}{f(\bar x)}.
\]
This follows by expanding the terms in which $b-1$ occurs exactly once.
We choose distinct nonzero elements $h_1,\ldots,h_p$ of the infinite field $k$, choose lifts $\widetilde h_i\in R$, and set $f_i=1+\widetilde h_iX$.
If $\sum_i c_i h_i\bar x/(1+h_i\bar x)=0$ with $c_i\in\Delta_0$, divide by $\bar x$ and clear denominators.
Then the polynomial $\sum_i c_i h_i\prod_{j\ne i}(1+h_jX)\in\Delta_0[X]$ has degree at most $p-1$ and vanishes at $\bar x$.
Since $[\Delta_0(\bar x):\Delta_0]=p$, this polynomial is zero.
Substitution of $X=-1/h_i$ gives $c_i=0$ for every $i$.
These logarithmic derivatives therefore form a basis of $\Delta_0(\bar x)$ over $\Delta_0$.
Finite products of the ratios thus realize every leading coefficient of a principal unit.

In the homogeneous-uniformizer type, we use $x^r$ as reference.
If $p\mid r$, principal units of the corner $C$ give all coefficients of that layer.
If $p\nmid r$, the coefficient of $1+a(bx)^r$ is $\bar a\bar b^r$.
Here $\Delta_0=\Delta$; choose $e_0$ with $\Delta^{p^{e_0}}\subset k$ and integers $s_0,t_0$ with $rs_0+p^{e_0}t_0=1$.
Every $z\in\Delta^\times$ can then be written $z=(z^{p^{e_0}})^{t_0}(z^{s_0})^r$.
We lift $\bar a=(z^{p^{e_0}})^{t_0}\in k^\times$ to the center $R$ and $\bar b=z^{s_0}\in\Delta^\times$ to $C^\times$.
The single principal unit $1+a(bx)^r$ realizes the leading coefficient $z$.
Hence every layer is realized in this type as well.

After finitely many layers, the norm of the deep remainder lies in $N_{T/F}V$.
Indeed, by separability of $T/F$ and nondegeneracy of the reduced trace, $\Tr_{T/F}\operatorname{Trd}_{B/T}$ is a nonzero $F$-linear functional, so after scaling by a power of the fixed uniformizer $\pi$ of $F$ we can choose $z\in J_C$ with $\Tr_{T/F}\operatorname{Trd}_{B/T}(z)\ne0$.
The polynomial
\[
 g(X)=N_{T/F}\Nrd_{B/T}(1+Xz)
\]
satisfies $g(0)=1$ and $g'(0)\ne0$, so by Hensel's lemma its image contains all sufficiently deep principal units of $F$.
Thus finitely many layer corrections and one deep norm correction suffice, without using an infinite product.
This proves the reduced-norm generation in both types.

Kato's norm-image formula and the purely inseparable residue extension give
\[
 \Nrd_{D/F}(U_D)=\ker\chi_{[D]}\cap U_F,
 \qquad V=\ker\chi_{[B]}\cap U_T.
\]
For example, if the reduced norm of a unit has residue one, that residue is the purely inseparable degree power of the residue of the unit, so the unit itself has residue one.
The projection formula and restriction formula for invariants give $V^G=\ker\chi_{p[D]}\cap U_F$.

Finally, Kato duality shows that the image of $\chi_{[D]}$ on $U_F$ has order $n$: the annihilator of $U_F$ is $\ker(\Br(F)\to\Br(F^{\mathrm{nr}}))$.
Indeed, the pairings with the two Witt components $\Br(k)$ and $H^1(k,\Q/\Z)$ give, respectively, the finite-order continuous characters of the valuation and residue-unit components of $F^\times/U_F$.
Ordinary local reciprocity and Kato duality show that these are all such characters.
Minimality of the index gives $m[D]\notin\ker(\Br(F)\to\Br(F^{\mathrm{nr}}))$, so the image on $U_F$ has order $n$ \cite[Theorem 2]{KatoExistence}.
We choose $u\in U_F$ such that $\chi_{[D]}(u)$ has order $n$.
Then $u^m=N_{F(x^p)/F}u$ belongs to $\ker\chi_{p[D]}\cap U_F$, and $\chi_{[D]}(u^m)$ has order $p$.
Since $\chi_{p[D]}=p\chi_{[D]}$ and $\chi_{[D]}(U_F)\subset\Q/\Z$ is cyclic of order $n$, the map $\chi_{[D]}$ identifies
\[
 (\ker\chi_{p[D]}\cap U_F)/(\ker\chi_{[D]}\cap U_F)
\]
with the subgroup of order $p$ of $\chi_{[D]}(U_F)$.
This quotient therefore has order $p$ and is generated by the class of $u^m$, so the group $\ker\chi_{[D]}\cap U_F$ and the element $u^m$ generate $V^G$.
\end{proof}

\begin{lemma}\label{lem:second-two-orbit-image}
In the situation of Lemma~\ref{lem:second-two-cyclic-cover}, assume that $K_2(B)=I_T(B)$, and let $Y=K(H,\bar H)$ be the relative spectrum of the cover $H=\mathcal Q(P)$ as above.
Then the map $\pi_2(Y_{hG})\to K_2(D)$, which forgets the cover through $K(B)_{hG}$ and passes to the generic point, satisfies
\[
 \im\bigl[\pi_2(Y_{hG})\longrightarrow K_2(D)\bigr]\subseteq I_F(D).
\]
\end{lemma}
\begin{proof}
We describe integral sources for the generators in Lemma~\ref{lem:review-homogeneous-unit-norm-generation}.
In the homogeneous-uniformizer type, use the relative Iwahori cover of $E_b$.
The relative corner of this cover is $\mathcal O_{F(w_b^p)}$, with degree-one group $U(F(w_b^p))$.
The element $w_b$ fixes $F(w_b^p)$, and $\operatorname{Ad}(b)$ acts trivially on $K_1(C)$; hence the class in $M$ of $a_b\in U(F(w_b^p))$ is $G$-invariant.
We take $a_b=N_{E_b/F(w_b^p)}z$ with $z\in U(E_b)$, so that by transitivity its reduced norm is $N_{F(w_b^p)/F}a_b=N_{E_b/F}z$.
The action and its carry are taken from the actual graded subring $\mathcal O_{E_b}=\bigoplus_{j=0}^{p-1}\mathcal O_{F(w_b^p)}w_b^j$, using $w_b^p$.

In the homogeneous-unit type, put $a_b=N_{E_b/F(w_b^p)}f(bx)$ and $a_1=N_{F(x)/F(x^p)}f(x)$.
Here $b\in1+J_C$, so $bx$ and $x$ have the same reduction, and the $p$-power norm of $\Delta/\Delta_0$ gives $\bar a_b=\bar a_1$.
Set $R_b=\mathcal O_{E_b}\times_\Delta\mathcal O_{F(x)}$ and $S_b=\tau_{\ge0}K(\mathcal Q(R_b))$.
Both projections are graded and have the same composite to the residue, so their difference, with its canonical residue nullhomotopy, gives a $G$-equivariant map $S_b\to Y$.
The corner unit pair $(a_b,a_1)$ is a $G$-invariant class of $\pi_1S_b$ whose image in $M$ is $[a_b]-[a_1]$.
Indeed, the degree-one unit $r=(bx,x)$ commutes with the degree-zero corner, and vertex translation fixes this unit pair.
The carry of the weak action is $r^p$.
For the additional class $u$, use the commutative relative cover of $F(x)$ and its corner principal unit $u\in U(F(x^p))$.

We apply Lemma~\ref{lem:review-generic-source-edge-generation} with these sources, $Y\to K(B)$, and $I_F(D)\subset K_2(D)$.
The map that forgets source orbits factors, as a spectrum map, through $K(E_b)$; at the generic point it is extension of scalars from $E_b$ to $D$, that is, the normalized transfer $\theta_{E_b}$ given by the $(E_b,D)$-bimodule $D$.
For the source $S_b$, it factors through the difference of the two maps at the generic point $R_b\otimes_RF=E_b\times F(x)$.
Thus the \emph{entire} image of the first filtration step lies in $I_F(D)$, through $\theta_{E_b}$ or $\theta_{E_b}-\theta_{F(x)}$, respectively.
The image of the corner $K_2(B)$ also lies in $I_F(D)$ by the inductive hypothesis $K_2(B)=I_T(B)$ and transitivity of transfer.
On the other hand, since $G$ is cyclic, we have natural isomorphisms $H_1(G,M)\simeq M^G/N_GM$ and $H_1(G,K_1(B))\simeq K_1(B)^G/N_GK_1(B)$.
Since $\SK_1(B)=0$, the reduced norm $\Nrd_{B/T}\colon K_1(B)\to\Nrd_{B/T}(B^\times)$ is a $G$-equivariant isomorphism, and $\Nrd_{B/T}(M)=V$ because the image of $U_C$ in $K_1(C)$ lies in $M$ and $\Nrd_{B/T}(M)\subseteq U_T\cap\ker\chi_{[B]}=V$.
Hence $\Nrd_{B/T}$ maps the image of $H_1(G,M)$ in $H_1(G,K_1(B))$ onto the image of $\Nrd_{B/T}(M^G)\subseteq V^G$ in $\Nrd_{B/T}(B^\times)^G/N_{T/F}\Nrd_{B/T}(B^\times)$.
By Lemma~\ref{lem:review-homogeneous-unit-norm-generation} and $N_{T/F}V=\Nrd_{B/T}(N_GM)$, the generic-point image of $H_1(G,M)$ is therefore generated by these source edges.
We therefore obtain the asserted inclusion in both types.
\end{proof}

\begin{proof}[Completion of the proof of Theorem~\ref{thm:second-two-wild}]
The case $n=1$ is split.
Assume $n>1$, and let $\xi\in\SK_2(D)$.
For $c=\partial_D\xi$, the boundary comparison in the proof of Lemma~\ref{lem:review-minimal-order-k1-support} gives $c=\partial_F\Nrd_2\xi=1$, so $\xi$ lifts to some $y\in K_2(P)$.
We first show that the support map $i_*\colon K_2(\Delta)\to K_2(P)$ vanishes.
Lifting Matsumoto symbols through a maximal subfield of ramification index one makes the reduction $r\colon K_2(P)\to K_2(\Delta)$ surjective.
For a uniformizer $x_0$ of $D$, put $\alpha=\operatorname{Ad}(x_0^{-1})_*$ on $K_2(P)$.
Additivity for the bimodule sequence $J_D\to P\to\Delta$ gives $i_*r=1-\alpha$, since $J_D=x_0P$ and the left action on $J_D$ is twisted by $\operatorname{Ad}(x_0^{-1})$, as in the proof of Lemma~\ref{lem:second-residue-zero-ell-integral}.
Conjugation by $x_0$ fixes the purely inseparable residue, so $r\alpha=r$ and $(1-\alpha)^2=i_*r(1-\alpha)=i_*r-i_*r\alpha=0$.
Since $x_0^n\in\pi P^\times$ and $\pi$ is central, $\operatorname{Ad}(x_0^n)$ is conjugation by a unit of $P$, so $1=\alpha^n=(1-(1-\alpha))^n=1-n(1-\alpha)$.
Thus $ni_*r=n(1-\alpha)=0$, and surjectivity of $r$ gives $ni_*=0$, so $i_*$ from the uniquely $p$-divisible group $K_2(\Delta)$ is zero.
By the localization sequence $K_2(\Delta)\xrightarrow{i_*}K_2(P)\to K_2(D)$, the map $K_2(P)\to K_2(D)$ is therefore injective.
Hence $n\xi=0$ gives $ny=0$, and $r(y)=0$ because $K_2(\Delta)$ is uniquely $p$-divisible.

Fix a finite separable splitting field $L_0/F$, and choose a depth $j$ at which Lemma~\ref{lem:app-congruence-separation} holds for $L_0$.
We lift the trace of $y$ by the surjection in degree two from absolute degree-zero $\mathrm{TC}$ to $\mathrm{TC}(P)^\wedge_p$ of Lemma~\ref{lem:second-two-degree-zero-surjective} and remove the residue component of the lift.
The order of the steps matters: in the homogeneous-unit type the residue orbit is removed by the $\mathrm{TC}$ class of the commutative subcover (in the homogeneous-uniformizer type, $\pi_2$ of the residue orbit already vanishes), the image with the cover forgotten is lifted to relative $\mathrm{TC}$ and compared there using $\pi_3\mathrm{TC}(\Delta)^\wedge_p=0$, and only then is it represented, at the prescribed depth $j$, by an integral class.
In the mixed-characteristic homogeneous-unit type, the image of the subtracted $\mathrm{TC}$ class with the cover forgotten is represented at depth $j$ by an integral class from $\mathcal O_E$, already compared on the residue and on the prescribed truncation.
By the equivalence in Lemma~\ref{lem:review-cyclic-zero-tc}, the remaining relative orbit class is $t_{\mathrm{rel}}\in\pi_2((\mathrm{TC}(H,\bar H)^\wedge_p)_{hG})$.
Here the outer $p$-completion in that lemma can be removed, since $\mathrm{TC}(H,\bar H)^\wedge_p$ is bounded below (the absolute terms lie in $\mathrm{Sp}_{\ge-1}$ and the relative term in $\mathrm{Sp}_{\ge-2}$) and $G$ is finite.
We then send $t_{\mathrm{rel}}$ to depth $j$: by the last assertion of Lemma~\ref{lem:app-orbit-stable-image}, its image at depth $j$ is represented by the trace of an integral element of $\pi_2(Y_{hG})$.
The generic-point image of that integral element lies in $I_F(D)$ by Lemma~\ref{lem:second-two-orbit-image}.
To compare with $y$, note that $r(y)=0$, so $y$ lifts to a relative class $y_{\mathrm{rel}}$.
In relative $\mathrm{TC}$ we have
\[
 \operatorname{tr}(y_{\mathrm{rel}})=f_*(t_{\mathrm{rel}})+b_{\mathrm{rel}},
\]
where $f_*$ forgets the cover and $b_{\mathrm{rel}}$ is the relative correction added back in the mixed-characteristic homogeneous-unit type (zero otherwise).
The absolute equality gives this relative one because $\pi_3\mathrm{TC}(\Delta)^\wedge_p=0$, so that the map from $\pi_2$ of relative $\mathrm{TC}$ to $\pi_2$ of absolute $\mathrm{TC}$ is injective.
We send this agreement to the prescribed finite truncation, where the relative trace comparison for nilpotent ideals (Lemma~\ref{lem:app-nilpotent-trace}) gives the agreement with integral coefficients.
Adding back the integral representatives, namely the representative from the cover and, in the mixed-characteristic homogeneous-unit type, the subtracted representative from $\mathcal O_E$, produces $z\in K_2(P)$ with $z|_D\in I_F(D)$ and $y-z\in\ker[K_2(P)\to K_2(P/\pi^jP)]$.
By Lemma~\ref{lem:app-congruence-separation}, we have $(y-z)|_D\in\theta_{L_0}K_2(L_0)$.
Thus $\xi\in I_F(D)$, and injectivity in Proposition~\ref{prop:review-splitting-transfer-norm-iso} gives $\xi=0$.
This closes the induction, and unramified descent in Lemma~\ref{lem:second-two-inertial-reduction} gives the result for the original $F,D$.
\end{proof}

We reduce Theorem~\ref{thm:second-two-vanishing} to Theorem~\ref{thm:second-two-wild}.

\begin{proof}[Proof of Theorem~\ref{thm:second-two-vanishing}]
Surjectivity is proved in Theorem~\ref{thm:second-two-surj}.
When the first residue field $k$ has characteristic zero, write $F=k((t))$ by Cohen's structure theorem, and let $B$ be the residue division algebra of $D$, with center $L$.
Since $L$ is a nonarchimedean local field, Theorem~\ref{thm:paper-local-norm} gives $\SK_1(B)=0$ and the isomorphism $\Nrd_{B/L,2}\colon K_2(B)\xrightarrow{\sim}K_2(L)$.
This isomorphism respects the action of $\Gal(L/k)$, so Proposition~\ref{prop:second-residue-zero} and cyclic Hilbert 90 for $K_2$ \cite[Theorem 14.1]{MerkurjevSuslin1982} give $\SK_2(D)=0$.
Suppose now that $k$ has characteristic $p>0$; then $k$ is a nonarchimedean local field.
For a prime $\ell\ne p$, an extension of degree prime to $\ell$ splitting the other primary Brauer components reduces the $\ell$-primary part of $\SK_2(D)$, by restriction and transfer, to that of an $\ell$-primary division algebra, and Proposition~\ref{prop:second-residue-zero-ell}(2) shows that it vanishes.
After an extension of degree prime to $p$ splitting the prime-to-$p$ Brauer components, restriction and transfer show that the $p$-primary part of $\SK_2$ is detected over that extension.
By Morita equivalence, the remaining cases are therefore exactly those of Theorem~\ref{thm:second-two-wild}.
\end{proof}
